% !TEX root = ../thesis.tex
\clearpage
\pagenumbering{Roman}
\section*{Zusammenfassung}
Hier könnte Ihre Zusammenfassung stehen.
\newpage
\section*{Abstract}
Under microwave irradiation, quasiparticle, Cooper pair, and multiple Andreev reflection (MAR) transport overlap within the same spectrum, and the voltage position of a spectral feature alone no longer identifies the process behind it. In this thesis, I separated these processes by combining independent static characterization with complete microwave amplitude maps. Static transport fixes the energy thresholds and the channel composition, the replica spacing sets the transferred charge, and the evolution with drive amplitude distinguishes incoherent tunneling from coherent phase locking. I applied this approach to aluminum junctions at millikelvin temperatures, from oxide tunnel barriers in the weak transmission limit to atomic contacts formed in a mechanically controllable break junction. An in situ calibration converts the microwave source output into the local ac voltage at the junction, so that every map can be compared with transport models on the physical drive coordinate.

In the first of two tunnel junctions, I reproduced the complete microwave driven evolution with the Tien-Gordon Model and resolved individual photon assisted quasiparticle replicas continuously to at least photon order 25. To my knowledge, no previous amplitude map of a superconducting tunnel junction resolves and continuously tracks individual replicas to a higher order by the same visual criterion. This result demonstrates quantitative photon assisted tunneling deep into the multiphoton regime and sets the reference for the more complex junctions that follow. In a second, nominally similar junction, I observed microwave induced rectification with a local current asymmetry reaching \qty{25}{\percent}. Dc bias, microwave amplitude, and frequency jointly control its magnitude and polarity, and the response lies beyond the bias symmetric Tien-Gordon description. Magnetic field and gate dependent control measurements support a microwave induced origin, and the narrow frequency interval in which the effect appears matches an electromagnetic mode of the sample environment identified independently. A resonantly driven Josephson weak link offers a plausible explanation for the drive threshold, the frequency sensitivity, and the polarity reversal, although the microscopic mechanism remains open.

For atomic contacts, channel resolved fits to the static MAR spectra yield the mesoscopic pincode, the set of individual channel transmissions. Followed along a complete opening trace, these transmissions evolve continuously within conductance plateaus and reorganize at abrupt jumps, revealing contact changes that stay hidden in the normal-state conductance alone. Using the pincode as independent input, I developed a deterministic threshold decomposition that partitions the static MAR conductance around its known thresholds and applies charge dependent photon dressing to each resulting order, yielding a threshold weighted description of photon assisted MAR (PAMAR). Combined with photon assisted incoherent Cooper pair tunneling (PAICPT), this framework reproduces the dominant finite-bias trajectories and their spectral weight evolution from moderate transmission to a nearly ballistic leading channel, without adjusting parameters to individual driven features. The static environmental analysis and the squared-Bessel amplitude evolution of the replicas originating from zero bias support their assignment to PAICPT rather than to a coherent Shapiro response. At the highest conductance investigated, the controlled few-channel description reaches its limit and an additional replica family appears at half the conventional Cooper pair spacing, a spacing at which, to my knowledge, PAICPT has not been reported before. Within the phenomenological description, this subharmonic PAICPT response matches the effective transfer of two Cooper pairs, although the measurements do not establish a simultaneous, correlated two-pair transfer or identify a unique microscopic mechanism.

Taken together, these measurements establish microwave amplitude spectroscopy as a charge sensitive method for separating overlapping transport processes by their static origin, transferred charge, and coherence. They connect quantitative photon assisted tunneling in the tunnel limit with a practical description of driven few-channel transport built from PAMAR and PAICPT, and they reveal two responses at the limits of that description, a tunable rectification tied to a mode of the sample environment and a subharmonic PAICPT response whose microscopic origin remains open.
